% GENERATED FILE. Edit canonical lessons, then run npm run generate:books.
% canonical-source-hash: fnv1a64:b78f7205726510de
% canonical-lessons: SA-C03-katham, SA-C03-bhavan-katham-asti, SA-S06-vowel-sign-aa, SA-C03-aham, SA-C03-kushalam, SA-S01-letter-ma, SA-W03-mama-guided-copy, SA-W03-mama-delayed-copy, SA-W03-mama-dictation, SA-C03-kushali-asmi, SA-W03-mama-nama-guided-copy, SA-W03-mama-nama-delayed-copy, SA-W03-mama-nama-dictation, SA-W03-namah-guided-copy, SA-W03-namah-delayed-copy, SA-W03-namah-dictation, SA-C03-na-cinta, SA-C03-practice, SA-S203-vowel-sign-u

\chapter{How Are You?}
\label{ch:responding}
\hlchaptermodality{3}

\begin{chapteropening}
\textbf{By the end of this chapter:} \emph{I can ask how someone is in Sanskrit, answer naturally, and return the question.}

Complete SA-C03-practice, the canonical closing checkpoint, and demonstrate the chapter promise: ask how someone is in Sanskrit, answer naturally, and return the question; then recall where katham, aham, kuśalam and cintā come from, why aham can drop, and pick out \sk{म} and \sk{◌ा}.
\end{chapteropening}

\section[\texorpdfstring{\sk{कथम्} (katham)}{कथम् (katham)}]{\sk{कथम्} (katham) — \textquotedblleft{}how,\textquotedblright{} another of the ka- questions}

\label{lesson:SA-C03-katham}



\begin{quote}
You met \emph{kim} (\textquotedblleft{}what\textquotedblright{}) in Chapter 2. Here is its sibling in the question-family.
\end{quote}

\begin{sounds}
\emph{ka · tham}. The \emph{th} is one sound — a \emph{t} with breath after it, never the \emph{th} of English \emph{thin}.
\end{sounds}

\begin{cousinweb}
\emph{katham} (\textquotedblleft{}how, in what way\textquotedblright{}) is from the interrogative stem \emph{ka-}, PIE \textbf{*k\textsuperscript{w}o-} — the same root that gave \emph{kim} and, in the west, English \textbf{wh-} and Latin \emph{qu-}. It heads the same \emph{ka-} question-family: \emph{kim} (what), \emph{kaḥ} (who), \emph{kadā} (when), \emph{kutra} (where), \emph{katham} (how). From \emph{katham} comes the abstract noun \emph{kathā} (\textquotedblleft{}a telling, a story\textquotedblright{} — \textquotedblleft{}the how of a thing\textquotedblright{}), and, worn down, Hindi's \emph{kaise}.
\end{cousinweb}

\subsection*{Your turn}
Say these aloud:

\begin{itemize}
\raggedright
  \item \textquotedblleft{}katham\textquotedblright{}
  \item the ka- question family — \emph{kim, kaḥ, kadā, kutra, katham}
\end{itemize}

\begin{tcolorbox}[breakable,colback=teal!4,colframe=teal!35!black,title={Before you move on}]
What root do Sanskrit's question-words share, and its western cousins? (The interrogative \emph{ka-}, PIE \emph{*k\textsuperscript{w}o-}; English \emph{wh-}, Latin \emph{qu-}.)
\end{tcolorbox}
\section[\texorpdfstring{\sk{भवान्} \sk{कथम्} \sk{अस्ति}? (bhavān katham asti?)}{भवान् कथम् अस्ति? (bhavān katham asti?)}]{\sk{भवान्} \sk{कथम्} \sk{अस्ति}? (bhavān katham asti?) — \textquotedblleft{}how are you?\textquotedblright{}}

\label{lesson:SA-C03-bhavan-katham-asti}



\begin{quote}
Three words you own, assembled into the first thing you ask after hello.
\end{quote}

\begin{cousinweb}
\emph{bhavān katham asti?} = \emph{bhavān} (respectful \textquotedblleft{}you\textquotedblright{}) + \emph{katham} (\textquotedblleft{}how\textquotedblright{}) + \emph{asti} (\textquotedblleft{}is\textquotedblright{}) — literally \textquotedblleft{}\textbf{your honour how is?}\textquotedblright{} Recall from Chapter 2 that \emph{bhavān} (\textquotedblleft{}your honour\textquotedblright{}) takes a \textbf{third-person} verb — so it is \emph{asti} (\textquotedblleft{}is\textquotedblright{}), not \textquotedblleft{}are.\textquotedblright{} The familiar version uses \emph{tvam} and the second-person verb: \emph{tvaṁ katham asi?} (\emph{asi}, \textquotedblleft{}you are\textquotedblright{}).
\end{cousinweb}

\begin{grammarlens}[title={Grammar Lens: the copula, three persons}]
You now hold the \textquotedblleft{}to be\textquotedblright{} verb across persons: \emph{asmi} (\textquotedblleft{}I am\textquotedblright{}), \emph{asi} (\textquotedblleft{}you are\textquotedblright{}), \emph{asti} (\textquotedblleft{}is\textquotedblright{}) — the ancestors, respectively, of English \textbf{am}, and (through Latin/Greek) of \emph{is/est}. \textquotedblleft{}How are you?\textquotedblright{} simply asks after a \emph{state}, and Sanskrit, unlike Tamil, has always kept a verb for it.
\end{grammarlens}

\subsection*{Your turn}
Say these aloud:

\begin{itemize}
\raggedright
  \item the respectful question — \emph{bhavān katham asti?}
  \item the familiar version — \emph{tvaṁ katham asi?}
  \item the copula trio — \emph{asmi} (am), \emph{asi} (are), \emph{asti} (is)
\end{itemize}

\begin{tcolorbox}[breakable,colback=teal!4,colframe=teal!35!black,title={Before you move on}]
Why is it \emph{asti} (\textquotedblleft{}is\textquotedblright{}) with \emph{bhavān}, and what are the three copula forms? (\emph{Bhavān} = \textquotedblleft{}your honour,\textquotedblright{} a 3rd-person subject; \emph{asmi/asi/asti} — am / you are / is.)
\end{tcolorbox}
\section[\texorpdfstring{\sk{ा} (ā)}{ा (ā)}]{\sk{◌ा} — one character, met inside words you already say}

\label{lesson:SA-S06-vowel-sign-aa}



\begin{quote}
Before the new one, the ones you have already met: \sk{आ} · \sk{भ} · \sk{◌ि}. Say what each of them does.

One character this time. One only — and you have been saying it for pages without knowing which mark on the page it was.
\end{quote}

\subsection*{Script you'll notice: \sk{◌ा}}
\textbf{\sk{◌ा}} — \emph{ā}.

It is a \textbf{vowel sign}. It is not a letter and never stands alone: it attaches to a consonant and \textbf{replaces} the \emph{a} built into it. Replaces, not adds — the \emph{a} is gone.

Where it sits: a vertical bar to the right of the consonant.

Worked through: \textbf{\sk{न}} + \textbf{\sk{◌ा}} = \textbf{\sk{ना}} — \emph{nā}.

You already say these words, and \sk{◌ा} is one of the shapes inside them — the rest of their shapes are still ahead of you, so here the character stands on its own:

\begin{itemize}
\raggedright
  \item \emph{namaskāraḥ} — hello (lit. \textquotedblleft{}the making of a bow\textquotedblright{})
  \item \emph{dhanyavādaḥ} — thank you
  \item \emph{svāgatam} — welcome (lit. \textquotedblleft{}well come\textquotedblright{})
  \item \emph{nāma} — name
\end{itemize}

\subsection*{Writing: \sk{◌ा} — follow the numbered strip}
\hlblockfigure{\detokenize{figures/SA-S06-vowel-sign-aa-filmstrip.pdf}}{How \sk{ा} is written, stroke by stroke}

Follow the numbered strip: start where movement 1 begins, and let each panel show you the next movement before your pen makes it. Then write \sk{◌ा} yourself — slowly, and larger than it is printed.

\begin{quote}
The strip shows \textbf{where to start the character and which way to travel}, in one attested order. That is taught with real variation from school to school, so the strip names its source beneath it: treat it as a sound way in, not the only one.
\end{quote}

\subsection*{Your turn}
\begin{itemize}
\raggedright
  \item \emph{Look:} at this, and find \sk{◌ा} in it
\end{itemize}

\begin{quote}
\sk{◌ा}
\end{quote}

\begin{itemize}
\raggedright
  \item \emph{Trace it:} \sk{◌ा} three times, saying \emph{ā} as you finish each one
  \item \emph{Look:} back at any page of this chapter and find \sk{◌ा} once more
\end{itemize}

\begin{tcolorbox}[breakable,colback=teal!4,colframe=teal!35!black,title={Before you move on}]
Which character is this — \sk{◌ा}? What vowel does it put on the consonant it attaches to — and what does it take off? (\textbf{Puts \emph{ā} on; takes the built-in \emph{a} off.}) Name one word you already say that contains it.
\end{tcolorbox}
\section[\texorpdfstring{\sk{अहम्} (aham)}{अहम् (aham)}]{\sk{अहम्} (aham) — \textquotedblleft{}I,\textquotedblright{} the ancestor of \emph{ego} and \emph{I}}

\label{lesson:SA-C03-aham}



\begin{quote}
To answer \textquotedblleft{}how are you?\textquotedblright{} you first need \textquotedblleft{}I\textquotedblright{} — and it is the grandparent of both English \emph{I} and the word \emph{ego}.
\end{quote}

\begin{sounds}
\emph{a · ham}. Two short beats, closing on \emph{m}.
\end{sounds}

\begin{cousinweb}
\emph{aham} (\textquotedblleft{}I\textquotedblright{}) is from Proto-Indo-European \textbf{*h\textsubscript{1}eǵ(H)om} — the direct ancestor of Latin \textbf{ego} (whence English \emph{ego}, \emph{egotist}), Greek \emph{egō}, and, worn all the way down, English \textbf{I} itself. So Sanskrit \emph{aham}, Latin \emph{ego}, and English \emph{I} are one word across three continents of time. Its possessive is the \emph{mama} (\textquotedblleft{}my\textquotedblright{}) from Chapter 2 — and, curiously, \textquotedblleft{}I\textquotedblright{} (\emph{aham}) and \textquotedblleft{}my\textquotedblright{} (\emph{mama}) come from \emph{different} ancient roots (\emph{*eǵ-} and \emph{*me-}), a quirk Sanskrit shares with English (\emph{I} vs. \emph{my/me}).
\end{cousinweb}

\begin{grammarlens}[title={Grammar Lens: \textquotedblleft{}I\textquotedblright{} can be dropped}]
Sanskrit verbs carry the person in their ending — \emph{asmi} already means \textquotedblleft{}\textbf{I} am\textquotedblright{} — so \emph{aham} is often left out. You still learn it first, because the answer to \emph{katham asi?} begins with it.
\end{grammarlens}

\subsection*{Your turn}
Say these aloud:

\begin{itemize}
\raggedright
  \item \textquotedblleft{}aham\textquotedblright{}
  \item the family — \emph{aham} / Latin \emph{ego} / English \textbf{I}
  \item the odd pair — \textquotedblleft{}I\textquotedblright{} (\emph{aham}, root \emph{*eǵ-}) and \textquotedblleft{}my\textquotedblright{} (\emph{mama}, root \emph{*me-}) differ, as in English \emph{I} vs. \emph{my}
\end{itemize}

\begin{tcolorbox}[breakable,colback=teal!4,colframe=teal!35!black,title={Before you move on}]
What PIE root is \emph{aham} from, and its cousins? (\emph{*h\textsubscript{1}eǵ(H)om}; Latin \emph{ego}, English \textbf{I}.)
\end{tcolorbox}
\section[\texorpdfstring{\sk{कुशलम्} (kuśalam)}{कुशलम् (kuśalam)}]{\sk{कुशलम्} (kuśalam) — \textquotedblleft{}well-being\textquotedblright{}}

\label{lesson:SA-C03-kushalam}



\begin{quote}
The word that answers \textquotedblleft{}how are you?\textquotedblright{} — one that once meant, of all things, \textquotedblleft{}skill in cutting sacred grass.\textquotedblright{}
\end{quote}

\begin{sounds}
\emph{ku · śa · lam}. The \emph{ś} is the \emph{sh} of \emph{ship}. Three short beats, closing on a bare \emph{m}.
\end{sounds}

\begin{cousinweb}
\emph{Kuśalam} (\textquotedblleft{}well-being, welfare, all-being-well\textquotedblright{}) is native Sanskrit, famously derived from \emph{kuśa} (the sacred grass) + a root of \textquotedblleft{}cutting\textquotedblright{} — so it first meant \textquotedblleft{}skilled {[}at gathering kuśa grass properly{]},\textquotedblright{} hence \textquotedblleft{}adept, fit, in good order,\textquotedblright{} hence \textquotedblleft{}well.\textquotedblright{} A whole idea of wellness grew from a ritual skill. Its adjective is \emph{kuśalī} (\textquotedblleft{}one who is well\textquotedblright{}).
\end{cousinweb}

\subsection*{Your turn}
Say these aloud:

\begin{itemize}
\raggedright
  \item well-being — \emph{kuśalam}
  \item its adjective (\emph{kuśalī}, \textquotedblleft{}one who is well\textquotedblright{})
  \item the surprising root (\emph{kuśa}, the sacred grass $\to$ \textquotedblleft{}skilled\textquotedblright{} $\to$ \textquotedblleft{}well\textquotedblright{})
\end{itemize}

\begin{tcolorbox}[breakable,colback=teal!4,colframe=teal!35!black,title={Before you move on}]
What does \emph{kuśalam} mean, and where does it come from? (\textquotedblleft{}Well-being\textquotedblright{}; from \emph{kuśa} grass $\to$ \textquotedblleft{}skilled\textquotedblright{} $\to$ \textquotedblleft{}well.\textquotedblright{})
\end{tcolorbox}
\section[\texorpdfstring{\sk{म} (ma)}{म (ma)}]{\sk{म} — one character, met inside words you already say}

\label{lesson:SA-S01-letter-ma}



\begin{quote}
Before the new one, the ones you have already met: \sk{भ} · \sk{◌ि} · \sk{◌ा}. Say what each of them does.

One character this time. One only — and you have been saying it for pages without knowing which mark on the page it was.
\end{quote}

\subsection*{Script you'll notice: \sk{म}}
\textbf{\sk{म}} — \emph{ma}.

It is a \textbf{consonant}, and in this script a consonant is never bare: it comes with an \emph{a} already in it. So it is not \emph{m}, it is \textbf{ma}.

What it is made of:

\begin{itemize}
\raggedright
  \item a top-to-bottom left stem joined to a clockwise lower loop and rightward crossbar
  \item a top-to-bottom right stem
  \item the top shirorekhā
\end{itemize}

You already say these, and every one of them has \sk{म} somewhere inside it:

\begin{itemize}
\raggedright
  \item \textbf{\sk{आम्} / \sk{न}} \emph{ām / na} — yes / no
  \item \textbf{\sk{नाम}} \emph{nāma} — name
  \item \textbf{\sk{मम}} \emph{mama} — my
\end{itemize}

\subsection*{Writing: \sk{म}}
\hlblockfigure{\detokenize{figures/SA-S01-letter-ma-filmstrip.pdf}}{How \sk{म} is written, stroke by stroke}

\begin{itemize}
\raggedright
  \item \textbf{1.} descend the left stem, curl left and clockwise around the lower loop, and continue right through the crossbar without lifting
  \item \textbf{2.} without lifting, climb up the right stem to the headline
  \item \textbf{3.} descend the right stem top-to-bottom
  \item \textbf{4.} lift and draw the top shirorekhā left-to-right
\end{itemize}

\textbf{Pen lifts: 1.} The pen comes up once.

\begin{quote}
The source shows three movements; most native writers draw \sk{म} in two strokes (80\% of HP Labs India's native-writer samples). This path keeps every movement in order and direction, and lifts only before the headline.
\end{quote}

\begin{quote}
This is one attested teaching order and not a national standard — handwriting here is taught with school-to-school variation. Source: JackPotte, ‘Devanagari m \sk{म}.gif’, strokes 1–3, Wikimedia Commons, 29 March 2009.
\end{quote}

\subsection*{Your turn}
\begin{itemize}
\raggedright
  \item \emph{Look:} at this, and find \sk{म} in it
\end{itemize}

\begin{quote}
\sk{आम्} / \sk{न}  ·  \sk{नाम}  ·  \sk{मम}
\end{quote}

\begin{itemize}
\raggedright
  \item \emph{Trace it:} \sk{म} three times, saying \emph{ma} as you finish each one
  \item \emph{Look:} back at any page of this chapter and find \sk{म} once more
\end{itemize}

\begin{tcolorbox}[breakable,colback=teal!4,colframe=teal!35!black,title={Before you move on}]
Which character is this — \sk{म}? What sound does it carry? (\emph{\textbf{ma}}.) Name one word you already say that contains it.
\end{tcolorbox}
\section[\texorpdfstring{\sk{मम} (mama)}{मम (mama)}]{\sk{मम} — one known letter becomes one known word}

\label{lesson:SA-W03-mama-guided-copy}



\begin{quote}
Say \textbf{\sk{मम}} — “my.” Point to its two equal parts: \textbf{\sk{म} + \sk{म}}. Nothing here is new vocabulary, and no new shape is hiding in the word.
\end{quote}

\subsection*{Writing — guided copy}
\hlblockfigure{\detokenize{figures/SA-W03-mama-guided-copy-filmstrip.pdf}}{How \sk{मम} is written, letter by letter, then one headline, stroke by stroke}

Keep \textbf{\sk{मम}} visible and follow the numbered strip: write the body of the first \textbf{\sk{म}}, then the body of the second, and only then draw one \textbf{śirorekhā} left to right across both, so the word hangs from one line. Most native writers leave the headline for last; some draw it first.

Read what you wrote: \textbf{\sk{मम}}, “my.” Two copies of \emph{ma} make \emph{mama}; neither consonant needs a separate vowel sign because each carries its short \emph{a}.

\begin{tcolorbox}[breakable,colback=teal!4,colframe=teal!35!black,title={Before you move on}]
How many written syllables did the word need? (\textbf{Two: \sk{म} + \sk{म}.}) What made the two letters look like one word? (\textbf{Their joined headline.})
\end{tcolorbox}
\section[\texorpdfstring{\sk{मम} (mama)}{मम (mama)}]{\sk{मम} — hide the word, then write it}

\label{lesson:SA-W03-mama-delayed-copy}



\begin{quote}
Look at \textbf{\sk{मम}} for five seconds. Find the two copies of \textbf{\sk{म}} and the headline that joins them. Then cover the word.
\end{quote}

\subsection*{Writing — delayed copy}
\hlblockfigure{\detokenize{figures/SA-W03-mama-delayed-copy-filmstrip.pdf}}{How \sk{मम} is written, letter by letter, then one headline, stroke by stroke}

Wait ten seconds. With the model still covered, write the word meaning “my.” There is no tracing guide and no romanized answer to copy.

Only after the word is complete, uncover \textbf{\sk{मम}}. Compare three things in order: two equal letter bodies, two right stems, and one joined headline. Repair only the first difference you find.

\begin{tcolorbox}[breakable,colback=teal!4,colframe=teal!35!black,title={Before you move on}]
Read the repaired word once and give its meaning. (\textbf{\sk{मम} — “my.”})
\end{tcolorbox}
\section[\texorpdfstring{\sk{मम} (mama)}{मम (mama)}]{One heard word — write “my”}

\label{lesson:SA-W03-mama-dictation}



\begin{quote}
Cover every answer below this line. You will hear one familiar Sanskrit word. There is no visible Devanagari model and no romanized answer.
\end{quote}

\subsection*{Writing — word dictation}
\emph{Hear:} the Sanskrit word meaning “my”

Write the word from sound and meaning alone. Stop after one attempt. Now uncover and compare: \textbf{\sk{मम}}. If needed, repair one letter body, one stem, or the joined headline — only the first difference.

\begin{tcolorbox}[breakable,colback=teal!4,colframe=teal!35!black,title={Before you move on}]
The hand has now moved from one dictated syllable, \textbf{\sk{न}}, to a whole known word, \textbf{\sk{मम}}. Say the word once as you read your answer.
\end{tcolorbox}
\section[\texorpdfstring{\sk{अहं} \sk{कुशली} \sk{अस्मि} (ahaṁ kuśalī asmi)}{अहं कुशली अस्मि (ahaṁ kuśalī asmi)}]{\sk{अहं} \sk{कुशली} \sk{अस्मि} (ahaṁ kuśalī asmi) — \textquotedblleft{}I am well\textquotedblright{}}

\label{lesson:SA-C03-kushali-asmi}



\begin{quote}
You have the question and you have the word. Now put the answer together.
\end{quote}

\begin{sounds}
\emph{a · haṁ · ku · śa · lī · as · mi}. The \emph{ṁ} is a hum through the nose that colours the vowel before it rather than a full \emph{m}; the \emph{ī} in the fifth beat is long.
\end{sounds}

\begin{cousinweb}
\emph{Ahaṁ kuśalī asmi} — \textquotedblleft{}\textbf{I am well}\textquotedblright{}: \emph{aham} (\textquotedblleft{}I\textquotedblright{}) + \emph{kuśalī} (\textquotedblleft{}well\textquotedblright{}) + \emph{asmi} (\textquotedblleft{}am\textquotedblright{}). Or simply \emph{kuśalam} (\textquotedblleft{}{[}all is{]} well\textquotedblright{}). The exchange:

\begin{quote}
— \emph{bhavān katham asti?} (\textquotedblleft{}How are you?\textquotedblright{}) — \emph{ahaṁ kuśalī asmi, dhanyavādaḥ.} (\textquotedblleft{}I am well, thank you.\textquotedblright{})
\end{quote}
\end{cousinweb}

\begin{grammarlens}[title={Grammar Lens: \emph{asmi} is the \textquotedblleft{}I\textquotedblright{} form of \textquotedblleft{}to be\textquotedblright{}}]
\emph{Asti} was \textquotedblleft{}is.\textquotedblright{} \emph{Asmi} is \textquotedblleft{}am\textquotedblright{} — the same verb wearing its first-person ending \emph{-mi}. That \emph{-mi} is the mark of \textquotedblleft{}I\textquotedblright{} all through the Sanskrit verb, and you will meet it again on every verb in the next chapter.
\end{grammarlens}

\subsection*{Your turn}
Say these aloud:

\begin{itemize}
\raggedright
  \item the question — \emph{bhavān katham asti?}
  \item the full reply — \emph{ahaṁ kuśalī asmi, dhanyavādaḥ}
  \item the ending that means \textquotedblleft{}I\textquotedblright{} (\emph{-mi}, as in \emph{asmi})
\end{itemize}

\begin{tcolorbox}[breakable,colback=teal!4,colframe=teal!35!black,title={Before you move on}]
Give the reply to \emph{bhavān katham asti?}, and say which ending marks \textquotedblleft{}I.\textquotedblright{} (\emph{Ahaṁ kuśalī asmi}; the ending \emph{-mi}.)
\end{tcolorbox}
\section[\texorpdfstring{\sk{मम} \sk{नाम} (mama nāma)}{मम नाम (mama nāma)}]{\sk{मम} \sk{नाम} — two known words, one visible phrase}

\label{lesson:SA-W03-mama-nama-guided-copy}



\begin{quote}
Read \textbf{\sk{मम} \sk{नाम}} — “my name.” Both words and every shape are already familiar. Leave one ordinary word space between them.
\end{quote}

\subsection*{Writing — guided phrase copy}
\hlblockfigure{\detokenize{figures/SA-W03-mama-nama-guided-copy-filmstrip.pdf}}{How \sk{मम} \sk{नाम} is written, word by word, each with its own headline, stroke by stroke}

Keep the phrase visible and follow the numbered strip, one word at a time.

\begin{enumerate}
\raggedright
  \item Write \textbf{\sk{मम}}: the body of each \textbf{\sk{म}}, then one headline across the word.
  \item Leave one word space.
  \item Write \textbf{\sk{नाम}}: \textbf{\sk{न}}, the stem of \textbf{\sk{◌ा}}, then \textbf{\sk{म}}, and last one headline across the word.
\end{enumerate}

The vowel sign changes \emph{na} to \emph{nā}; it does not begin a separate letter. Most native writers leave the headline for last; some draw it first. Read the finished phrase once: \textbf{\sk{मम} \sk{नाम}}, “my name.”

\begin{tcolorbox}[breakable,colback=teal!4,colframe=teal!35!black,title={Before you move on}]
Which sign makes the long vowel in \textbf{\sk{नाम}}? (\textbf{\sk{◌ा}}, attached to \textbf{\sk{न}}.) Where does the pen leave a visible gap? (\textbf{Between the two words.})
\end{tcolorbox}
\section[\texorpdfstring{\sk{मम} \sk{नाम} (mama nāma)}{मम नाम (mama nāma)}]{\sk{मम} \sk{नाम} — hide the phrase, keep its two-word shape}

\label{lesson:SA-W03-mama-nama-delayed-copy}



\begin{quote}
Look at \textbf{\sk{मम} \sk{नाम}} for five seconds. Point to the word space and to \textbf{\sk{◌ा}}. Then cover the model.
\end{quote}

\subsection*{Writing — delayed phrase copy}
\hlblockfigure{\detokenize{figures/SA-W03-mama-nama-delayed-copy-filmstrip.pdf}}{How \sk{मम} \sk{नाम} is written, word by word, each with its own headline, stroke by stroke}

Wait ten seconds. Write “my name” in Sanskrit with no visible answer and no romanization. Keep the two words apart, but join the headline within each word.

Only after the phrase is complete, uncover \textbf{\sk{मम} \sk{नाम}}. Compare meaning order, word space, letters, and vowel sign in that order. Repair only the first difference.

\begin{tcolorbox}[breakable,colback=teal!4,colframe=teal!35!black,title={Before you move on}]
Read the phrase once. Which word came first? (\textbf{\sk{मम}}, “my.”)
\end{tcolorbox}
\section[\texorpdfstring{\sk{मम} \sk{नाम} (mama nāma)}{मम नाम (mama nāma)}]{Two heard words — write “my name”}

\label{lesson:SA-W03-mama-nama-dictation}



\begin{quote}
Cover every answer below this line. You will hear a familiar two-word fragment. There is no visible Devanagari model and no romanized answer.
\end{quote}

\subsection*{Writing — phrase dictation}
\emph{Hear:} the Sanskrit phrase meaning “my name”

Write the two words from sound and meaning alone. Stop after one attempt. Now uncover and compare: \textbf{\sk{मम} \sk{नाम}}. Check meaning order, the one word space, and the long-vowel sign in \textbf{\sk{नाम}}. Repair only the first difference.

\begin{tcolorbox}[breakable,colback=teal!4,colframe=teal!35!black,title={Before you move on}]
The runway now reaches a two-word phrase. It does not yet reach a full introduction: \textbf{\sk{अस्ति}}, names, punctuation, and connected text remain later rungs, after their shapes are taught and recalled in order.
\end{tcolorbox}
\section[\texorpdfstring{\sk{नमः} (namaḥ)}{नमः (namaḥ)}]{\sk{नमः} — the known bow, now ending in a visible breath}

\label{lesson:SA-W03-namah-guided-copy}



\begin{quote}
Say the first half of \emph{namaste}: \emph{namas}, “a bow, a salutation.” That meaning is already yours. Said by itself at a pause, its final \emph{s} appears as the visarga: \emph{namaḥ}. This is the same known word, not new vocabulary.
\end{quote}

\subsection*{Writing — guided visarga copy}
Keep \textbf{\sk{नमः}} visible and copy it once.

\begin{enumerate}
\raggedright
  \item Write \textbf{\sk{न}}, then \textbf{\sk{म}}, joining their headline into one word.
  \item Put \textbf{\sk{ः}} to the right of the final syllable. Its two dots do not join the headline.
  \item Read the result aloud: \emph{na-ma-ḥ} — “a bow, a salutation.”
\end{enumerate}

The page now shows all three things the ear hears: \textbf{\sk{न} + \sk{म} + \sk{ः}}.

\begin{tcolorbox}[breakable,colback=teal!4,colframe=teal!35!black,title={Before you move on}]
Which mark carries the final breath? (\textbf{\sk{ः}}, the visarga.) Is \emph{namaḥ} a new word here? (\textbf{No. It is the known \emph{namas} from \emph{namaste}, said at a pause.})
\end{tcolorbox}
\section[\texorpdfstring{\sk{नमः} (namaḥ)}{नमः (namaḥ)}]{\sk{नमः} — hide the word, keep the final breath}

\label{lesson:SA-W03-namah-delayed-copy}



\begin{quote}
Look at \textbf{\sk{नमः}} for five seconds. Point to \textbf{\sk{न}}, \textbf{\sk{म}}, and the visarga \textbf{\sk{ः}}. Say \emph{namaḥ}, “a bow, a salutation.” Then cover the model.
\end{quote}

\subsection*{Writing — delayed visarga copy}
Wait ten seconds. Write the known word with no visible answer and no romanization. Keep \textbf{\sk{नम}} beneath one headline and place the two visarga dots to its right.

Only after the word is complete, uncover \textbf{\sk{नमः}}. Compare the two letters, their order, and the final mark. Repair only the first difference.

\begin{tcolorbox}[breakable,colback=teal!4,colframe=teal!35!black,title={Before you move on}]
Read what you wrote. Which known piece of \emph{namaste} did you recover? (\emph{\textbf{Namas}}, “a bow,” in its pause form \emph{namaḥ}.)
\end{tcolorbox}
\section[\texorpdfstring{\sk{नमः} (namaḥ)}{नमः (namaḥ)}]{A heard bow — write the visarga from sound and meaning}

\label{lesson:SA-W03-namah-dictation}



\begin{quote}
Cover every answer below this line. You will hear the familiar word meaning “a bow, a salutation,” the first half of \emph{namaste}, now said by itself. There is no visible Devanagari model and no romanized answer.
\end{quote}

\subsection*{Writing — visarga dictation}
\emph{Hear:} \emph{namaḥ} — “a bow, a salutation”

Write the word from sound and meaning alone. Stop after one attempt. Now uncover and compare: \textbf{\sk{नमः}}. Check \textbf{\sk{न}}, \textbf{\sk{म}}, and the final visarga in that order. Repair only the first difference.

\begin{tcolorbox}[breakable,colback=teal!4,colframe=teal!35!black,title={Before you move on}]
The runway now carries a Sanskrit-specific mark from recognition to independent transcription. Which later debts remain? (\textbf{Vocalic ṛ, conjuncts, connected text, and independent recall across a longer passage.})
\end{tcolorbox}
\section[\texorpdfstring{\sk{न} \sk{चिन्ता} (na cintā)}{न चिन्ता (na cintā)}]{\sk{न} \sk{चिन्ता} (na cintā) — \textquotedblleft{}no worry\textquotedblright{}}

\label{lesson:SA-C03-na-cinta}



\begin{quote}
When someone thanks you, this gentle reply says, simply, \textquotedblleft{}don't give it a thought.\textquotedblright{}
\end{quote}

\begin{sounds}
\emph{na cin · tā}. The \emph{c} is the \emph{ch} of \emph{church}, never a \emph{k}; the final \emph{ā} is long.
\end{sounds}

\begin{cousinweb}
\emph{Na cintā} literally means \textquotedblleft{}\textbf{{[}there is{]} no worry / no anxiety}\textquotedblright{} — used for \textquotedblleft{}never mind / no trouble / you're welcome.\textquotedblright{} Two atoms: \emph{na} (\textquotedblleft{}not\textquotedblright{}), the ancient negative from Chapter 1 — PIE \textbf{*ne}, behind English \textbf{no}, \textbf{not}, Latin \emph{nōn}, and every Indo-Aryan \textquotedblleft{}no.\textquotedblright{} And \emph{cintā} (\textquotedblleft{}anxious thought\textquotedblright{}), from the root \emph{cint} (\textquotedblleft{}to think, ponder\textquotedblright{}) — the same \emph{cintā} borrowed into Hindi and Urdu (\emph{cintā mat karo}, \textquotedblleft{}don't worry\textquotedblright{}). So the reply pairs the family's oldest negative with a verb of thinking: \textquotedblleft{}think nothing of it.\textquotedblright{}
\end{cousinweb}

\subsection*{Your turn}
Say these aloud:

\begin{itemize}
\raggedright
  \item \textquotedblleft{}na cintā\textquotedblright{}
  \item its two atoms (\emph{na} \textquotedblleft{}not\textquotedblright{} + \emph{cintā} \textquotedblleft{}worry\textquotedblright{})
  \item the exchange — \emph{dhanyavādaḥ} / \emph{na cintā}
\end{itemize}

\begin{tcolorbox}[breakable,colback=teal!4,colframe=teal!35!black,title={Before you move on}]
What does \emph{na cintā} literally say, and what is its negative \emph{na} cousin to? (\textquotedblleft{}No worry\textquotedblright{}; PIE \emph{*ne} — English \emph{no/not}, Latin \emph{nōn}.)
\end{tcolorbox}
\section[Practice]{Chapter 3 — The how-are-you exchange}

\label{lesson:SA-C03-practice}



\begin{quote}
No new words. The whole responding cycle, from atoms you now own.
\end{quote}

\subsection*{The exchange}
Say these aloud:

\begin{itemize}
\raggedright
  \item \textquotedblleft{}How are you?\textquotedblright{} (respectful) — \emph{bhavān katham asti?}
  \item \textquotedblleft{}I am well, thank you.\textquotedblright{} — \emph{ahaṁ kuśalī asmi, dhanyavādaḥ.}
  \item \textquotedblleft{}No worry / you're welcome.\textquotedblright{} — \emph{na cintā.}
  \item the familiar question — \emph{tvaṁ katham asi?}
\end{itemize}

\subsection*{How to answer: the atoms}
Say these aloud:

\begin{itemize}
\raggedright
  \item the ka- question-word for \textquotedblleft{}how\textquotedblright{} (\emph{katham})
  \item \textquotedblleft{}I\textquotedblright{} and the copula trio (\emph{aham}; \emph{asmi/asi/asti})
\end{itemize}

\subsection*{Your turn — where the words come from}
Say these aloud:

\begin{itemize}
\raggedright
  \item the stem \emph{katham} shares with \emph{kim}, and its western cousins (\emph{ka-}, PIE \emph{*k\textsuperscript{w}o-}; English \emph{wh-}, Latin \emph{qu-})
  \item the cousins of \emph{aham} (Latin \emph{ego}, English \textbf{I})
  \item why \emph{aham} is often left out, and the ending that marks \textquotedblleft{}I\textquotedblright{} (\emph{asmi} already means \textquotedblleft{}I am\textquotedblright{}; \emph{-mi})
  \item the surprising source of \emph{kuśalam} (\emph{kuśa}, the sacred grass $\to$ \textquotedblleft{}skilled\textquotedblright{} $\to$ \textquotedblleft{}well\textquotedblright{})
  \item what \emph{cintā} is, and its root (anxious thought; \emph{cint}, \textquotedblleft{}to think\textquotedblright{})
\end{itemize}

\subsection*{Script — the chapter's two shapes}
\begin{quote}
\sk{◌ा} · \sk{म}
\end{quote}

\begin{itemize}
\raggedright
  \item \emph{Point to:} \emph{ma}, then the \emph{ā}-sign
  \item \emph{Say it:} what each does, and a word with both (\textbf{\sk{म}}, \emph{ma}, its \emph{a} built in; \textbf{\sk{◌ा}}, a bar to the right that puts \emph{ā} on and takes the \emph{a} off; \emph{nāma})
\end{itemize}

\begin{grammarlens}[title={What you've built: the roots you now carry}]
\begin{itemize}
\raggedright
  \item \emph{katham} $\leftarrow$ the interrogative \emph{ka-} (PIE \emph{*k\textsuperscript{w}o-}) $\to$ English \emph{how/what}, Latin \emph{qu-}
  \item \emph{aham} $\leftarrow$ \emph{*h\textsubscript{1}eǵ(H)om} $\to$ Latin \textbf{ego}, English \textbf{I}
  \item \emph{asmi} $\leftarrow$ \emph{*h\textsubscript{1}es-} $\to$ English \textbf{am}; \emph{asti} $\to$ \emph{is}, \emph{hai}
  \item \emph{na} (in \emph{na cintā}) $\leftarrow$ \emph{*ne} $\to$ English \emph{no/not}, Latin \emph{nōn}
\end{itemize}
\end{grammarlens}

\begin{tcolorbox}[breakable,colback=teal!4,colframe=teal!35!black,title={Before you move on}]
A respected stranger asks \emph{bhavān katham asti?} — give a full reply, and answer their thanks. (\emph{Ahaṁ kuśalī asmi, dhanyavādaḥ} — and to thanks, \emph{na cintā}.)
\end{tcolorbox}
\section[\texorpdfstring{\sk{ु} (u)}{ु (u)}]{\sk{◌ु} — one character, met inside words you already say}

\label{lesson:SA-S203-vowel-sign-u}



\begin{quote}
Before the new one, the ones you have already met: \sk{◌ि} · \sk{◌ा} · \sk{म}. Say what each of them does.

One character this time. One only — and you have been saying it for pages without knowing which mark on the page it was.
\end{quote}

\subsection*{Script you'll notice: \sk{◌ु}}
\textbf{\sk{◌ु}} — \emph{u}.

It is a \textbf{vowel sign}. It is not a letter and never stands alone: it attaches to a consonant and \textbf{replaces} the \emph{a} built into it. Replaces, not adds — the \emph{a} is gone.

Where it sits: a small hook below the consonant.

Worked through: \textbf{\sk{न}} + \textbf{\sk{◌ु}} = \textbf{\sk{नु}} — \emph{nu}.

You already say this word, and \sk{◌ु} is one of the shapes inside it — the rest of its shapes are still ahead of you, so here the character stands on its own:

\begin{itemize}
\raggedright
  \item \emph{kuśalam} — well-being
\end{itemize}

\subsection*{Writing: \sk{◌ु} — follow the numbered strip}
\hlblockfigure{\detokenize{figures/SA-S203-vowel-sign-u-filmstrip.pdf}}{How \sk{ु} is written, stroke by stroke}

Follow the numbered strip: start where movement 1 begins, and let each panel show you the next movement before your pen makes it. Then write \sk{◌ु} yourself — slowly, and larger than it is printed.

\begin{quote}
The strip shows \textbf{where to start the character and which way to travel}, in one attested order. That is taught with real variation from school to school, so the strip names its source beneath it: treat it as a sound way in, not the only one.
\end{quote}

\subsection*{Your turn}
\begin{itemize}
\raggedright
  \item \emph{Look:} at this, and find \sk{◌ु} in it
\end{itemize}

\begin{quote}
\sk{◌ु}
\end{quote}

\begin{itemize}
\raggedright
  \item \emph{Trace it:} \sk{◌ु} three times, saying \emph{u} as you finish each one
  \item \emph{Look:} back at any page of this chapter and find \sk{◌ु} once more
\end{itemize}

\begin{tcolorbox}[breakable,colback=teal!4,colframe=teal!35!black,title={Before you move on}]
Which character is this — \sk{◌ु}? What vowel does it put on the consonant it attaches to — and what does it take off? (\textbf{Puts \emph{u} on; takes the built-in \emph{a} off.}) Name one word you already say that contains it.
\end{tcolorbox}
